\documentclass[11pt,a4paper]{article}
\usepackage{turkos}

\begin{document}
\phaseheader{4}{Hardware Interrupts: PIC, Timer and Keyboard}{Let devices interrupt the CPU, keep time with the PIT, read the keyboard, and give Turk-OS its first interactive prompt}{22}{28}

\ataglance{2 weeks}{3}{Phase 3: GDT, IDT, exception stubs, \code{register_interrupt_handler}}{A remapped 8259 PIC, IRQ
dispatch with end-of-interrupt handling, a 100\,Hz timer with uptime and \code{sleep_ms}, a PS/2 keyboard driver
(shift, caps lock, optionally the Turkish Q layout), and a kernel monitor prompt with commands like \code{help},
\code{uptime}, \code{mem} and \code{reboot}.}

\section{Why this phase}

So far the kernel only does what \code{kmain} tells it, in order. Real systems are driven by events: a key press, a
timer tick, a disk finishing a read. Devices announce these events with \textbf{interrupts}, and in this phase you
connect the first two devices to the IDT you built in Phase 3.

The timer matters more than it looks. A periodic tick is how the kernel gets control back from a running program,
which makes preemptive multitasking possible in Phase 8, and it is also the kernel's clock for timeouts and
\code{sleep}. The keyboard gives you input, and together with the screen that makes Turk-OS interactive for the first
time. You will end the phase with a small command monitor inside the kernel. It is not a real shell, since that has
to be a user program (Phase 14), but you will use it constantly as a test console from now on.

\section{Concepts you need}

\subsection{The 8259 Programmable Interrupt Controller}

Devices don't wire straight into the CPU. On a PC they raise an \textbf{IRQ} line on an interrupt controller, the
classic 8259A PIC, which forwards one interrupt at a time to the CPU along with a vector number. There are two PICs in
a cascade: the master handles IRQ 0--7, and the slave handles IRQ 8--15 and reports through the master's IRQ 2
(Figure~\ref{fig:pic}).

\begin{figure}[htbp]
\centering
\begin{tikzpicture}[node distance=0.4cm and 1.2cm, every node/.style={font=\sffamily\scriptsize}]
  \node[tkhi, minimum width=2.2cm, minimum height=1.4cm] (cpu) {\textbf{CPU}\\INTR pin};
  \node[tkbox, minimum width=2.6cm, minimum height=2.0cm, left=2.7cm of cpu] (m) {\textbf{Master PIC}\\ports \code{0x20}/\code{0x21}\\vectors 32--39};
  \node[tkbox, minimum width=2.6cm, minimum height=2.0cm, below=0.9cm of m] (s) {\textbf{Slave PIC}\\ports \code{0xA0}/\code{0xA1}\\vectors 40--47};
  \draw[tkarrow] (m) -- node[above, tknote] {interrupt + vector} (cpu);
  \draw[tkarrow] (s.north) -- node[right, tknote] {cascade on IRQ 2} (m.south);
  \node[tknote, align=right, anchor=east] (mi) at ([xshift=-0.3cm]m.west) {IRQ 0 timer (PIT)\\IRQ 1 keyboard\\IRQ 2 cascade\\IRQ 3--4 serial\\IRQ 7 printer / spurious};
  \node[tknote, align=right, anchor=east] (si) at ([xshift=-0.3cm]s.west) {IRQ 8 real-time clock\\IRQ 12 PS/2 mouse\\IRQ 14 primary ATA disk\\IRQ 15 secondary ATA / spurious};
  \draw[tkarrow] (mi.east) -- (m.west);
  \draw[tkarrow] (si.east) -- (s.west);
\end{tikzpicture}
\caption{The two cascaded PICs after Turk-OS remaps them. IRQ 14 returns in Phase 12.}
\label{fig:pic}
\end{figure}

Three things about the PIC trip everyone up. First, the BIOS programs the master to use vectors 8--15, which collide
with CPU exceptions: the timer's IRQ 0 would arrive as vector 8 and look exactly like a double fault. So the very
first job is to \textbf{remap} the PICs to vectors 32--47. Second, after handling an IRQ you must send an
\textbf{end-of-interrupt} (EOI) command, or the PIC never delivers that IRQ again; for IRQs 8--15, both the slave
and the master need one. Third, each PIC has a \textbf{mask} register where a set bit disables an IRQ line. You will
unmask only the devices you have drivers for.

\subsection{Timekeeping with the PIT}

The 8253/8254 Programmable Interval Timer counts down from a value you choose, driven by a fixed
1\,193\,182\,Hz oscillator, and raises IRQ 0 each time the count reaches zero. For a 100\,Hz tick the divisor is
$1\,193\,182 / 100 \approx 11\,932$. Each tick the kernel increments a counter, and everything else follows from it:
uptime, \code{sleep}, timeouts, and in Phase 8 the scheduler's time slices. \MOS\,\S5.5 separates \emph{clock
hardware} from \emph{clock software} the same way, and also describes soft timers, a useful idea for later.

\subsection{The PS/2 keyboard}

A keyboard controller (the 8042) sits between the keyboard and the CPU. When you press or release a key it puts a
\textbf{scan code} in its data port \code{0x60} and raises IRQ 1. Scan codes identify physical keys, not characters.
In the default scan code set 1, pressing a key sends its \emph{make code} and releasing it sends the same code with
bit 7 set (the \emph{break code}). Some keys send an \code{0xE0} prefix first. Turning key events into characters is
your job, and it needs a little state: which modifiers (shift, ctrl, alt) are held and whether caps lock is on. The
last step is a \textbf{layout} table from key to character, and a Turkish Q keyboard needs a different table from a US
one. The controller will not report another key until you have read port \code{0x60}.

\subsection{Doing as little as possible in an interrupt}

An interrupt handler runs at an unpredictable moment, with interrupts disabled by the interrupt gate, on top of
whatever the CPU was doing. It should therefore do the minimum and get out: read the hardware, store the data, send
the EOI. The slow work (translating keys, echoing, parsing commands) happens outside, in normal code that takes data
from a buffer. \MOS\,\S5.3.1 and \OSDI\,\S3.2.2 describe this split, and Figure~\ref{fig:key} shows it for a key
press. The ring buffer between the two halves is your first shared data structure between interrupt context and
normal code. Phase 9 comes back to making that safe in general.

\begin{figure}[htbp]
\centering
\begin{tikzpicture}[node distance=0.35cm and 0.42cm, every node/.style={font=\sffamily\scriptsize}]
  \node[tkbox, text width=1.1cm] (k) {key press};
  \node[tkbox, text width=1.6cm, right=of k] (c) {8042 puts scan code in \code{0x60}};
  \node[tkbox, text width=1.5cm, right=of c] (p) {PIC: IRQ 1 $\rightarrow$ vector 33};
  \node[tkbox, text width=1.9cm, right=of p] (s) {stub + \code{isr_handler} + EOI};
  \node[tkhi, text width=2.8cm, right=of s] (h) {\code{keyboard_handler}:\\read port, store};
  \node[tkok, text width=1.5cm, below=0.7cm of h] (rb) {ring buffer};
  \node[tkbox, text width=2.6cm, left=0.9cm of rb] (m) {monitor: \code{kbd_getchar()}, echo, parse};
  \draw[tkarrow] (k) -- (c); \draw[tkarrow] (c) -- (p); \draw[tkarrow] (p) -- (s); \draw[tkarrow] (s) -- (h);
  \draw[tkarrow] (h) -- (rb); \draw[tkarrow] (rb) -- (m);
  \node[tknote, above=0.1cm of s] {interrupt context};
  \node[tknote, below=0.1cm of m] {normal kernel code};
\end{tikzpicture}
\caption{A key press, from the keyboard to the monitor. Only the red box runs with interrupts off.}
\label{fig:key}
\end{figure}

\section{Reading plan}

\begin{readingplan}
\must & \LOB & Ch.~6, rest: \emph{Programmable Interrupt Controller (PIC)}, \emph{Reading Input from the Keyboard} & 43--45 & Tasks 4.1, 4.2 and 4.4. \\
\must & \LOB & Ch.~14, section \emph{Programmable Interval Timer} & 74 & Task 4.3. \\
\must & \MOS & \S5.5 \emph{Clocks}: clock hardware, clock software, soft timers & 388--394 & What a tick is for, and what to build on it. \\
\must & \OSTEP & Ch.~36 \S36.3--36.4 (the canonical protocol; lowering CPU overhead with interrupts) & 407--409 & Polling versus interrupts, now with real devices. \\
\should & \MOS & \S5.2 \emph{Principles of I/O Software}; \S5.3.1 \emph{Interrupt Handlers} & 351--357 & Programmed versus interrupt-driven I/O; what a handler must do. \\
\should & \MOS & \S5.6.1 \emph{Input Software} (keyboard software) & 394--399 & Scan codes, raw versus cooked input, line editing. \\
\should & \OSDI & \S2.8 \emph{The Clock Task in MINIX 3} & 204--214 & A real clock driver: tick handling, alarms, accounting. \\
\should & \OSDI & \S3.8.5 \emph{Implementation of the Keyboard Driver} & 350--357 & Scan code processing and keymaps in MINIX. \\
\could & \OSDI & \S3.1.4 \emph{Interrupts}; \S3.2.2 \emph{Interrupt Handlers} & 226--227, 231 & Short, clear summary of the hardware side. \\
\could & \OSDI & \S3.8.1--3.8.2 \emph{Terminal Hardware}, \emph{Terminal Software} & 303--316 & Canonical mode and special keys, preview of Phase 14. \\
\end{readingplan}

\begin{minixbox}
MINIX programs the PIC in \texttt{kernel/i8259.c (08100)} with the same ICW1--ICW4 sequence as Task~4.1. Its clock
interrupt handler in \texttt{kernel/clock.c (10400)} does almost nothing: it updates the time, charges the tick to
the running process, and only wakes the clock task if something (an expired alarm or time slice) needs attention.
That is the ``do little in the handler'' rule in a real kernel. The keyboard driver in
\texttt{drivers/tty/keyboard.c (15200)} buffers raw scan codes in its interrupt handler and translates them later
through a loadable keymap, which is how MINIX supports many layouts.
\end{minixbox}

\section{Build it}

\task{Remap and mask the PIC}
Send the four initialisation words to both PICs, moving IRQ 0--7 to vectors 32--39 and IRQ 8--15 to 40--47, then mask
everything except the timer and keyboard. Give each write a moment with \code{io_wait()}; old hardware needed it and
it costs nothing.

\begin{lst}{c}{kernel/arch/i386/pic.c}
#define PIC1_CMD  0x20
#define PIC1_DATA 0x21
#define PIC2_CMD  0xA0
#define PIC2_DATA 0xA1

void pic_remap(void)
{
    outb(PIC1_CMD,  0x11); io_wait();   /* ICW1: begin initialisation, ICW4 follows */
    outb(PIC2_CMD,  0x11); io_wait();
    outb(PIC1_DATA, 0x20); io_wait();   /* ICW2: master vector offset = 32          */
    outb(PIC2_DATA, 0x28); io_wait();   /* ICW2: slave vector offset  = 40          */
    outb(PIC1_DATA, 0x04); io_wait();   /* ICW3: a slave is attached to IRQ 2       */
    outb(PIC2_DATA, 0x02); io_wait();   /* ICW3: slave's cascade identity is 2      */
    outb(PIC1_DATA, 0x01); io_wait();   /* ICW4: 8086/88 mode                       */
    outb(PIC2_DATA, 0x01); io_wait();
    outb(PIC1_DATA, 0xFC);              /* mask: only IRQ 0 and IRQ 1 enabled       */
    outb(PIC2_DATA, 0xFF);              /* mask every slave IRQ for now             */
}

void pic_send_eoi(uint8_t irq)
{
    if (irq >= 8)
        outb(PIC2_CMD, 0x20);           /* slave first for IRQ 8..15 */
    outb(PIC1_CMD, 0x20);
}
\end{lst}

\task{IRQ stubs and dispatch}
Add sixteen stubs for vectors 32--47 to \code{isr.s}. They push a dummy error code and the vector, then reuse
\code{isr_common}, so IRQs get the same \code{struct regs} as exceptions. Install them in \code{idt_init} with type
\code{0x8E}. In \code{isr_handler}, route vectors 32--47 to the registered handler and then send the EOI.

\begin{lst}{asm}{kernel/arch/i386/isr.s (additions)}
%macro IRQ 2
irq%1:
    push dword 0            ; IRQs have no error code
    push dword %2           ; vector number = 32 + IRQ
    jmp  isr_common
%endmacro

IRQ 0, 32
IRQ 1, 33
; ... one line per IRQ ...
IRQ 15, 47
\end{lst}

\begin{lst}{c}{kernel/arch/i386/isr.c (addition at the top of isr_handler)}
if (r->int_no >= 32 && r->int_no < 48) {
    if (handlers[r->int_no])
        handlers[r->int_no](r);
    pic_send_eoi(r->int_no - 32);
    return;
}
\end{lst}

\task{The timer}
Program PIT channel 0 for 100\,Hz and count ticks in a handler. Build \code{timer_ticks()},
\code{uptime_ms()} and \code{sleep_ms()} on top. Keep the counter a \code{volatile uint32_t}: on a 32-bit CPU, a
64-bit counter is read in two halves, and a tick can land between them and give a torn value. At 100\,Hz a 32-bit
counter lasts about 497 days.

\begin{lst}{c}{kernel/drivers/pit.c}
#define PIT_BASE_HZ 1193182

static volatile uint32_t ticks;

static void timer_handler(struct regs *r) { (void)r; ticks++; }

void pit_init(uint32_t hz)
{
    uint16_t divisor = PIT_BASE_HZ / hz;
    outb(0x43, 0x36);                   /* channel 0, low byte then high byte, mode 3 */
    outb(0x40, divisor & 0xFF);
    outb(0x40, divisor >> 8);
    register_interrupt_handler(32, timer_handler);
}

void sleep_ms(uint32_t ms)
{
    uint32_t end = ticks + ms / 10;     /* 10 ms per tick at 100 Hz */
    while ((int32_t)(end - ticks) > 0)  /* this comparison survives wrap-around */
        __asm__ volatile ("hlt");       /* sleep until the next interrupt */
}
\end{lst}

Enable interrupts only after everything they might touch is ready. The order in \code{kmain} is:
\code{gdt_init()}, \code{idt_init()}, \code{pic_remap()}, \code{pit_init(100)}, \code{keyboard_init()}, then
\code{__asm__ volatile ("sti")}.
\verify{Print the tick count once a second for a minute; it goes up by 100 each time. Compare \code{sleep_ms(5000)}
with a stopwatch.}

\task{The keyboard driver}
Write \code{kernel/drivers/keyboard.c}. The handler reads the scan code, updates modifier state on make and break
codes, and for ordinary key presses translates through a layout table and puts the character in a ring buffer (your
Phase 0 code). Start with the US layout; the table below covers the main block of keys in scan code set 1.

\begin{lst}{c}{kernel/drivers/keyboard.c}
static const char keymap_us[128] = {
    0, 27, '1','2','3','4','5','6','7','8','9','0','-','=','\b',   /* 0x00-0x0E */
    '\t','q','w','e','r','t','y','u','i','o','p','[',']','\n',      /* 0x0F-0x1C */
    0,  'a','s','d','f','g','h','j','k','l',';','\'','`',           /* 0x1D-0x29 */
    0,  '\\','z','x','c','v','b','n','m',',','.','/', 0, '*',       /* 0x2A-0x37 */
    0,  ' ',                                                        /* 0x38-0x39 */
};
/* keymap_us_shift[] is the same table with the shifted characters. */

static bool shift, ctrl, caps;
static struct ringbuf kbd_buf;

static void keyboard_handler(struct regs *r)
{
    (void)r;
    uint8_t sc = inb(0x60);             /* always read, or no further keys arrive */
    bool released = sc & 0x80;
    sc &= 0x7F;

    switch (sc) {
    case 0x2A: case 0x36: shift = !released; return;   /* left / right shift */
    case 0x1D:            ctrl  = !released; return;   /* left ctrl          */
    case 0x3A: if (!released) caps = !caps;  return;   /* caps lock toggles  */
    }
    if (released)
        return;

    char c = keymap_us[sc];
    bool letter = (c >= 'a' && c <= 'z');
    if (shift ^ (caps && letter))
        c = keymap_us_shift[sc];
    if (c)
        rb_put(&kbd_buf, c);
}
\end{lst}

The consumer side waits for a key with the processor halted. Checking the buffer and then halting has a tiny race:
if the key arrives between the check and \code{hlt}, you sleep until the next tick. On x86 the classic cure is to
disable interrupts for the check and use \code{sti; hlt} together, because \code{sti} only takes effect after the
following instruction, so no interrupt can slip in between them.

\begin{lst}{c}{kernel/drivers/keyboard.c (consumer)}
char kbd_getchar(void)
{
    char c;
    for (;;) {
        __asm__ volatile ("cli");
        if (rb_get(&kbd_buf, &c)) {
            __asm__ volatile ("sti");
            return c;
        }
        __asm__ volatile ("sti; hlt");  /* enable interrupts and sleep, atomically */
    }
}
\end{lst}
\verify{Letters, digits, punctuation, shift and caps lock all produce the right characters. Holding a key produces
auto-repeat, because the keyboard itself resends the make code.}

\task{The Turkish Q layout (optional)}
Add a second table and a way to switch between them, for example a \code{layout tr} command. On a Turkish Q
keyboard the letter keys sit on the same scan codes as their US neighbours: \code{0x17} gives ı (and I with shift),
\code{0x1A} ğ, \code{0x1B} ü, \code{0x27} ş, \code{0x28} i (and İ with shift), \code{0x33} ö and \code{0x34} ç.
Because several of these characters are outside ASCII, let the keyboard produce Unicode code points instead of
\code{char}, and let the console convert them for the screen. Code page 437 can show ç, ö and ü directly; print a
placeholder for ğ, ş and ı until Phase 15 gives you a font that has them.

\task{Line input and the kernel monitor}
Write \code{readline(buf, size)}, which echoes characters, handles backspace (it uses your VGA driver's
\texttt{\textbackslash b}) and returns on Enter. Then write the monitor loop: print a prompt, read a line, split it at
spaces, and look the first word up in a command table. Every later phase adds commands here.

\begin{center}
\small\renewcommand{\arraystretch}{1.25}
\begin{tabularx}{\linewidth}{@{}l X@{}}
\toprule
\tkth{Command} & \tkth{What it does} \\
\midrule
\code{help} & list all commands with one-line descriptions (from the command table itself) \\
\code{clear}, \code{echo <text>}, \code{color <fg> <bg>} & exercise the VGA driver \\
\code{uptime}, \code{ticks} & show time since boot from the PIT counter \\
\code{mem} & memory size from the Multiboot information (a real report arrives in Phase 5) \\
\code{crash <kind>} & run one of your Phase 3 exception tests \\
\code{reboot} & wait until \code{inb(0x64) & 0x02} is clear, then \code{outb(0x64, 0xFE)} to pulse the CPU reset line \\
\bottomrule
\end{tabularx}
\end{center}
\verify{A session like \code{turk> uptime} prints the right time; backspace edits the line; \code{reboot} restarts
QEMU and you see the boot banner again.}

\section{Testing and debugging}

Run with \code{-d int -no-reboot} once to watch the interrupts arrive: you should see a stream of \code{v=20} (timer)
entries and a \code{v=21} for every key event. The QEMU monitor command \code{info pic} shows the PIC state,
including the mask and in-service registers.

\begin{pitfalls}
Exactly one timer interrupt ever arrives & No EOI sent, so the PIC never delivers IRQ 0 again & Send the EOI after every IRQ, including when no handler is registered. \\
Double fault (vector 8) as soon as you run \code{sti} & PIC not remapped: IRQ 0 arrives as vector 8 & Remap before \code{sti}. This is the classic first-IRQ bug. \\
The keyboard works for one key, then stops & Port \code{0x60} not read in the handler, or no EOI & Read the scan code unconditionally at the start of the handler. \\
Timer much too fast or slow & Divisor wrong, or overflowed 16 bits & For 100\,Hz it is 11\,931; the minimum frequency is about 19\,Hz. \\
Random characters appear or disappear & The handler and the reader race on the ring buffer indices & One producer and one consumer with \code{volatile} indices is safe here; Phase 9 generalises it. \\
Occasional interrupt on vector 39 (IRQ 7) or 47 & Spurious IRQ from the PIC & Read the in-service register (OCW3 \code{0x0B}); if the bit is clear, it was spurious: skip the EOI (for IRQ 15, send it to the master only). \\
\end{pitfalls}

\section{Milestone}

\begin{milestonebox}[Turk-OS is interactive]
\begin{checklist}
  \item PIC remapped to 32--47; only the timer and keyboard are unmasked; every IRQ gets an EOI.
  \item Ticks arrive at 100\,Hz; \code{uptime} is accurate and \code{sleep_ms} works.
  \item Typing echoes correctly with shift, caps lock and backspace.
  \item The monitor runs \code{help}, \code{clear}, \code{uptime}, \code{mem}, \code{crash} and \code{reboot}.
  \item Left idle for ten minutes at the prompt, the kernel neither crashes nor spins: it halts between interrupts.
\end{checklist}
\end{milestonebox}

\begin{questionbox}
\begin{enumerate}
  \item Why must the PICs be remapped, and what exactly goes wrong if you skip it?
  \item What is an EOI, and why do IRQs 8--15 need two of them?
  \item Compute the PIT divisor for 1000\,Hz. What would that frequency cost and what would it buy you?
  \item How is a key release distinguished from a key press in scan code set 1?
  \item Why should an interrupt handler do as little as possible? Where does the rest of the work go?
  \item Why is \code{sti; hlt} safe while \code{sti} followed later by \code{hlt} is not?
\end{enumerate}
\end{questionbox}

\section{Going further}

Read the real-time clock (CMOS ports \code{0x70}/\code{0x71}) and add a \code{date} command that prints the date in
Turkish format (\code{25.09.2026 14:05}). Make the PC speaker beep using PIT channel 2. Add arrow-key history to
\code{readline} (arrows send \code{0xE0}-prefixed codes). If you want to prepare for multiprocessor support, read
about the local APIC and I/O APIC, which replace the PIC on modern machines; you will need them only if you take the
SMP track in Phase 15.

\nextphase{5}{Physical Memory Management}
\booklegend
\end{document}
